\chapter{Introduction}
\label{chap:intro}

Database management systems are critical infrastructure in modern software. SQLite is the most widely deployed database engine in the world, embedded in every major smartphone operating system, web browser, and the Python standard library, with an estimated three trillion active installations \citep{sqlite}. Despite extensive testing and code review, memory safety vulnerabilities continue to be discovered in SQLite at a steady pace, affecting billions of downstream devices and applications. The persistence of such defects in a mature, heavily audited codebase motivates the need for systematic automated vulnerability discovery.

Fuzzing is one of the most effective automated techniques for finding software vulnerabilities \citep{fuzzingsurvey}. Coverage-guided greybox fuzzers such as AFL \citep{afl} and libFuzzer \citep{libfuzzer} use edge coverage feedback to steer input generation toward unexplored program regions. However, byte-level mutation fuzzers fail on structured inputs like SQL: random mutations produce parser-rejected inputs that never reach deep program logic where vulnerabilities reside. Grammar-based fuzzing addresses this by generating inputs from a context-free grammar, guaranteeing syntactic validity. Nautilus \citep{nautilus} combines grammar-based generation with coverage feedback, and several tools target databases specifically: Squirrel \citep{squirrel}, SQLsmith \citep{sqlsmith}, SQLRight \citep{sqlright}, and Griffin \citep{griffin}. However, none of these systems pay enough attention to the grammar itself --- they use general-purpose SQL grammars that treat all constructs equally, with no domain knowledge about which structural patterns are most likely to trigger vulnerabilities.

This thesis presents DBMS-Nautilus, a grammar-based greybox fuzzing system for SQLite that focuses on grammar engineering as the primary lever for vulnerability discovery. The core idea is a \textit{structural primitives} grammar methodology: production rules encode SQL attack patterns associated with vulnerability classes --- schema setup with generated columns, stress queries with window functions, boundary function calls with integer overflow values --- without embedding any proof-of-concept exploit input. The grammar contains 526 production rules, iteratively refined based on evidence from coverage profiling and crash analysis.

In summary, the contributions of this thesis are:
\begin{itemize}
    \item A \textit{structural primitives} grammar methodology for SQL fuzzing with zero proof-of-concept contamination, where production rules encode vulnerability-triggering patterns as composable non-terminals rather than fixed exploit strings.
    \item A context-free grammar consisting of 526 production rules organized into four structural pattern categories --- schema topology, query stress, boundary values, and validation operations --- each targeting distinct SQLite subsystems where vulnerabilities cluster, iteratively refined based on CVE root-cause analysis and experimental feedback.
    \item An empirical evaluation across four SQLite versions: with seed rules, DBMS-Nautilus rediscovers 4 of 4 target CVEs through compositional generation. Without seed rules, it discovers 13 distinct bug classes across 48 unique stack hashes --- including 4 independent CVE rediscoveries, 6 non-CVE bugs that crash production builds across up to 10 consecutive releases, and 3 sanitizer-only detections. A comparative evaluation against a generic EBNF baseline shows $108\times$ more unique root causes per campaign.
\end{itemize}

The remainder of this thesis is organized as follows. Chapter~\ref{chap:background} presents background on fuzzing, grammar-based generation, DBMS testing tools, and sanitizer oracles. Chapter~\ref{chap:method} presents the DBMS-Nautilus design: system architecture, the structural primitives grammar methodology organized by pattern categories, and the cross-pollination property that enables compositional bug discovery. Chapter~\ref{chap:experiments} reports the experimental evaluation structured around three research questions: whether DBMS-Nautilus can rediscover known CVEs, whether it can find new bugs, and how it compares to a generic EBNF baseline grammar in terms of throughput, coverage, and crash diversity. The final chapter concludes with limitations and future work.
